% TOP_PROBLEM: {"id": 3270, "title": "Hoàng-Reed Conjecture", "category_id": 3, "category": {"id": 3, "name": "graph_theory", "display_name": "Graph Theory", "description": "Problems involving graphs, networks, and their properties.", "slug": "graph-theory", "order_index": 3, "created_at": "2026-07-31T15:26:25.671Z"}, "status": "open", "problem_number": "OPG-47282", "set_id": 12}
% FIRST_POSTED: 2026-10-02
% ATTEMPT_STATUS: unresolved
% UnsolvedMath ID 3270; source catalogue code OPG-47282
% !TeX program = lualatex
% GPT-6 Astra Ultra research attempt; independently checked partial work, 2026-10-02.
\documentclass[11pt,a4paper]{article}
\usepackage[margin=25mm]{geometry}
\usepackage{fontspec}
\setmainfont{lmroman10-regular.otf}[BoldFont=lmroman10-bold.otf,
 ItalicFont=lmroman10-italic.otf,BoldItalicFont=lmroman10-bolditalic.otf]
\usepackage{amsmath,amssymb,mathtools}
\usepackage{unicode-math}
\setmathfont{latinmodern-math.otf}
\AtBeginDocument{\let\setminus\smallsetminus}
\usepackage{xurl,hyperref}
\hypersetup{colorlinks=true,urlcolor=blue}
\setcounter{secnumdepth}{0}
\setlength{\parindent}{0pt}
\setlength{\parskip}{5pt}
\setlength{\emergencystretch}{3em}
\begin{document}
\section*{Hoàng-Reed Conjecture}\label{3270}
{\small UnsolvedMath ID 3270 · OPG-47282 · Graph Theory · OpenGarden}
\subsection{Definitions and mathematical statement}
Let $D=(V,A)$ be a nonempty finite loopless simple digraph, allowing opposite
arcs. Write $d^+(v)$ for the number of distinct vertices $w$ with
$vw\in A$. A directed cycle is a subdigraph following a cyclic sequence
of distinct vertices; cycles of length two are allowed when both
opposite arcs occur.

For each integer $k\ge1$, the Hoàng--Reed conjecture asks whether the
condition $d^+(v)\ge k$ for every vertex guarantees directed cycles
$C_1,\ldots,C_k$ which can be ordered so that
\[
 \left|V(C_j)\cap\bigcup_{i=1}^{j-1}V(C_i)\right|\le1
                         \qquad(2\le j\le k).
\]
Thus each newly chosen cycle is either disjoint from all earlier cycles
or meets their entire union at just one vertex. In particular the
cycles are arc-disjoint. They are not required to be vertex-disjoint,
and an arbitrarily large number of them may use one common vertex.

The condition is stronger than saying that every pair of cycles shares
at most one vertex: a new cycle cannot meet two earlier cycles at two
different vertices. The existence of a suitable ordering is part of
the conclusion. Extra arcs of $D$ between vertices of the chosen cycles
are ignored, since the union is not required to be induced.

For $k=1$ the assertion reduces to finding a directed cycle in a finite
digraph with no sink. The general question seeks a structured collection
of $k$ cycles from a local outdegree assumption, without any minimum
indegree or connectivity requirement.
\subsection{Short English statement}
Does minimum outdegree k force k directed cycles that can be added one by one, each meeting all previous cycles in at most one vertex?
\subsection{Research attempt: arc packing, overlap obstruction, and a girth reduction}
\paragraph{Scope and literature check.}
GPT-6 Astra Ultra, 2 October 2026. The intersection with the entire
previous union, rather than each previous cycle separately, is the
key quantifier. Welhan [S3] proved the Ho\`ang--Reed conjecture for
minimum outdegree three, so older source discussions must not be
read as a current claim that this case is unknown. The argument
below is an independent elementary analysis of weaker packings and
symmetric digraphs; it does not reprove that stronger published case.
It also verifies exactly how the conjecture would imply the usual
short-cycle bound.

\paragraph{1. Minimum outdegree always yields $k$ arc-disjoint cycles.}
Every finite digraph with positive minimum outdegree contains a
cycle: follow successive outgoing arcs until a vertex repeats, and
extract a simple cycle. Delete its arcs. At any vertex, this deletion
removes either one outgoing arc or none, so the remaining minimum
outdegree is at least $k-1$. Repeating the operation $k$ times gives
$k$ arc-disjoint directed cycles. Before step $i$, where counting
starts at one, the remaining minimum outdegree is at least $k-i+1$,
so each step is justified.

This proof controls arcs used at a vertex, but does not control how
many vertices a later cycle shares with earlier cycles. Its output
is therefore weaker than the required ordered intersection property.
No minimum indegree hypothesis is needed for this weaker conclusion.

\paragraph{2. Pairwise small intersections do not repair the argument.}
Consider the three directed triangles
\[
 (a,x,b,a),\qquad (b,y,c,b),\qquad (c,z,a,c).
\]
All six displayed vertices are distinct except for the repetitions
explicitly indicating the triangle ends. The first and second
triangles intersect at $b$, the second and third at $c$, and the
first and third at $a$. They are arc-disjoint and every pair shares
only one vertex. Nevertheless, in any ordering the last triangle
meets the union of the first two in two distinct vertices. Thus
this particular collection has no admissible ordering.

The example is an obstruction to upgrading a packing certificate,
not a counterexample to the conjecture: the digraph consisting of
these arcs does not satisfy minimum outdegree three. Keeping that
hypothesis distinction is essential. It shows that even proving
pairwise intersections at most one would leave a further step.

\paragraph{3. The symmetric subclass is immediate for every $k$.}
Suppose every arc has its opposite and minimum outdegree is at least
$k$. Choose a vertex $v$ and distinct outneighbours
$w_1,\ldots,w_k$. The $k$ digons
\[
              v\to w_i\to v\qquad(1\le i\le k)
\]
are distinct and arc-disjoint. Each new one meets the earlier union
exactly in $v$, so every ordering works. This proves the full desired
intersection pattern on all symmetric digraphs, without a regularity
or connectivity assumption. It relies on length-two cycles being
allowed in the canonical convention; it would not apply to oriented
graphs, where opposite arcs are absent.

\paragraph{4. Exact implication for the shortest cycle.}
Let $g$ denote directed girth, the length of a shortest directed
cycle. If the conjectured ordered collection exists, its first
cycle has at least $g$ vertices, and every subsequent cycle
contributes at least $g-1$ previously unused vertices. Therefore
\[
 n\ge g+(k-1)(g-1)=k(g-1)+1.
\]
Rearranging the integer inequality gives
\[
 g\le\left\lfloor\frac{n-1}{k}\right\rfloor+1
       =\left\lceil\frac nk\right\rceil.
\]
This is precisely the Caccetta--H\"aggkvist short-cycle conclusion.
The ceiling identity holds also when $k$ divides $n$, as follows
by writing $n=qk+r$ with $0\le r<k$. Arc-disjointness alone would
not give the vertex count; the stronger ordered overlap is what
makes the implication work.

\paragraph{Verification and first remaining gap.}
The script [S9] checks all six orders of the three-triangle
example, constructs the symmetric digon certificates for
$1\le k\le8$, and checks the displayed ceiling identity on bounded
integer inputs. These computations accompany exact arguments.
Deleting whole cycle vertex sets could enforce disjointness but may
remove many outgoing neighbours at once. Deleting only cycle arcs
preserves the inductive degree bound but loses the intersection
condition. No intermediate deletion or exchange procedure overcoming
both obstacles has been established here.

\paragraph{Final status.}
UNRESOLVED for arbitrary digraphs. The symmetric subclass, the weaker
arc-packing bound, and the precise short-cycle implication are proved.
\subsection{Sources}
\begin{itemize}
\item[S1] ULAM.AI and the UnsolvedMath Contributors, supplied problems.json, record 3270 (OPG-47282), OpenGarden collection. Catalogue identity and source transcription; CC BY 4.0.
\url{https://huggingface.co/datasets/ulamai/UnsolvedMath}.
\item[S2] Open Problem Garden, Hoàng-Reed Conjecture. Original problem and discussion; the mathematical formulation below is expanded from the supplied source record.
\url{http://www.openproblemgarden.org/op/hoand_reed_conjecture}.
\item[S3] Manuel Welhan, \emph{The Ho\`ang--Reed Conjecture for $\delta^+=3$}, Discrete Mathematics 310 (2010), 1932--1939. \url{https://doi.org/10.1016/j.disc.2010.03.001}.
\item[S9] GPT-6 Astra Ultra, exact finite checks accompanying this attempt (2 October 2026). Repository script, Python standard library only. \texttt{scripts/check\_attempt\_batch03\_digraphs\_2026\_10\_02.py}.
\end{itemize}
\end{document}
